\documentclass[11pt,a4paper]{article}
\usepackage[utf8]{inputenc}
\usepackage[T1]{fontenc}
\usepackage{geometry}
\geometry{margin=1in}
\usepackage{hyperref}
\usepackage{enumitem}
\usepackage{booktabs}
\usepackage{longtable}
\usepackage{xcolor}
\usepackage{titlesec}
\usepackage{amsmath,amssymb}

\titleformat{\section}{\Large\bfseries\color{blue!70!black}}{}{0em}{}
\titleformat{\subsection}{\large\bfseries\color{blue!50!black}}{}{0em}{}
\titleformat{\subsubsection}{\normalsize\bfseries\color{blue!30!black}}{}{0em}{}

\title{\textbf{Replication Plan}\\[0.5em]
\large Integer Sequences from Configurations in the Hausdorff Metric Geometry via Edge Covers of Bipartite Graphs\\[0.3em]
\normalsize OSTI ID: 1997354}
\author{Auto-generated Replication Blueprint}
\date{\today}

\begin{document}
\maketitle
\tableofcontents
\newpage

%---------------------------------------------------------------------
\section{Paper Overview}
%---------------------------------------------------------------------
This paper derives integer sequences arising from counting configurations in finite metric spaces under the Hausdorff metric. The key mathematical objects are bipartite graphs connecting points of two finite configurations, where edges represent distance relationships. The paper enumerates edge covers of these bipartite graphs to count the number of distinct Hausdorff distance configurations. All computations are purely combinatorial/algebraic and produce verifiable integer sequences (many registered in the OEIS).

%---------------------------------------------------------------------
\section{Detailed Workflow}
%---------------------------------------------------------------------

\subsection{Phase 1: Mathematical Framework}
\begin{enumerate}[label=\textbf{1.\arabic*}]
  \item \textbf{Implement Hausdorff distance computation.}
    \begin{equation*}
      d_H(A, B) = \max\left(\max_{a \in A} \min_{b \in B} d(a,b),\; \max_{b \in B} \min_{a \in A} d(a,b)\right)
    \end{equation*}
  \item \textbf{Enumerate finite configurations.}
    \begin{itemize}
      \item For a discrete metric space $X = \{1, 2, \ldots, n\}$, enumerate all non-empty subsets $A, B \subseteq X$.
      \item Compute $d_H(A, B)$ for all pairs.
    \end{itemize}
\end{enumerate}

\subsection{Phase 2: Bipartite Graph Construction}
\begin{enumerate}[label=\textbf{2.\arabic*}]
  \item \textbf{Construct the bipartite graph $G(A, B)$.}
    \begin{itemize}
      \item Vertices: elements of $A$ on one side, elements of $B$ on the other.
      \item Edges: connect $a \in A$ to $b \in B$ if $d(a, b) \leq r$ for a given radius $r$.
    \end{itemize}
  \item \textbf{Enumerate edge covers.}
    \begin{itemize}
      \item An edge cover is a set of edges such that every vertex is incident to at least one edge.
      \item Count edge covers using inclusion-exclusion or the permanent-like formula from the paper.
    \end{itemize}
\end{enumerate}

\subsection{Phase 3: Integer Sequence Generation}
\begin{enumerate}[label=\textbf{3.\arabic*}]
  \item \textbf{Systematically vary parameters} ($n$, $|A|$, $|B|$, distance threshold).
  \item \textbf{Compute the counting sequences} for each parameter combination.
  \item \textbf{Tabulate results} in the format matching the paper's tables.
  \item \textbf{Cross-reference with OEIS} to verify known sequences.
\end{enumerate}

\subsection{Phase 4: Closed-Form Verification}
\begin{enumerate}[label=\textbf{4.\arabic*}]
  \item \textbf{Verify closed-form expressions} given in the paper against brute-force enumeration.
  \item \textbf{Verify recurrence relations} for generating larger terms.
  \item \textbf{Test generating functions} where provided.
\end{enumerate}

\subsection{Phase 5: Validation}
\begin{enumerate}[label=\textbf{5.\arabic*}]
  \item Compare all generated sequences term-by-term with the paper's stated values.
  \item Verify OEIS entries match (by A-number).
  \item Extend sequences beyond what's in the paper as an additional check.
\end{enumerate}

%---------------------------------------------------------------------
\section{Datasets}
%---------------------------------------------------------------------

\begin{longtable}{p{4.5cm} p{3.5cm} p{6cm}}
\toprule
\textbf{Dataset / Source} & \textbf{Type} & \textbf{Location / Notes} \\
\midrule
\endfirsthead
\toprule
\textbf{Dataset / Source} & \textbf{Type} & \textbf{Location / Notes} \\
\midrule
\endhead
OEIS (Online Encyclopedia of Integer Sequences) & Reference integer sequences & \url{https://oeis.org/} \\
\midrule
Paper's tables of sequences & Reference & Extracted from the paper \\
\midrule
All other data & Synthetic (enumeration) & Generated by the algorithms \\
\bottomrule
\end{longtable}

%---------------------------------------------------------------------
\section{Models, Tools, and Software}
%---------------------------------------------------------------------

\begin{longtable}{p{4.5cm} p{2.5cm} p{6.5cm}}
\toprule
\textbf{Tool / Library} & \textbf{Version} & \textbf{Purpose / URL} \\
\midrule
\endfirsthead
\toprule
\textbf{Tool / Library} & \textbf{Version} & \textbf{Purpose / URL} \\
\midrule
\endhead
Python & $\geq 3.8$ & Primary implementation language \\
\midrule
SageMath & $\geq 9.5$ & Combinatorial enumeration, graph theory \newline \url{https://www.sagemath.org/} \\
\midrule
NetworkX & $\geq 2.6$ & Bipartite graph construction and analysis \newline \url{https://networkx.org/} \\
\midrule
SymPy & $\geq 1.10$ & Symbolic computation, generating functions \newline \url{https://www.sympy.org/} \\
\midrule
itertools (stdlib) & built-in & Combinatorial enumeration of subsets \\
\midrule
NumPy & $\geq 1.21$ & Array operations \\
\midrule
Matplotlib & $\geq 3.4$ & Visualization of graph structures \\
\bottomrule
\end{longtable}

\subsection{Hardware Requirements}
\begin{itemize}
  \item Any modern laptop. Enumeration is exponential in $n$ but the paper works with small $n$.
  \item For extending sequences to larger $n$: more RAM and multi-core CPU helpful.
\end{itemize}

\end{document}
